\documentclass[a4paper,12pt,french]{report}

% =======================
% IMPORTATION DES PACKAGES
% =======================

\usepackage[T1]{fontenc}
\usepackage[utf8]{inputenc}
\usepackage[margin=2cm]{geometry}
\usepackage{theorem}
\usepackage{lmodern}
\usepackage{amssymb, amsmath, amsfonts}
\usepackage{babel}

\renewcommand{\familydefault}{\sfdefault}

\pagestyle{headings}

\theoremstyle{break}
\theoremheaderfont{\scshape}
\theorembodyfont{\upshape}

\newtheorem{defin}{Définition}[section]
\newtheorem{theo}[defin]{Théorème}
\newtheorem{prop}[defin]{Proposition}

\title{MÉSURE ET INTÉGRATION}
\date{\today}
\author{}

\begin{document}
\maketitle
\textsc{Précision} \\
Ceci est juste des notes pour m'aider à comprendre la théorie de la mesure et intégration tout en apprenant latex.
Mais si cela peut vous aider, vous pouvez le consulter et le télécharger.

\tableofcontents

\chapter{THÉORIE DES ENSEMBLES}

\section{Théorie des ensembles}
\begin{defin}
Soit $X$ un ensemble non vide et $(A_{i})_{i \in \mathbb{N}}.$ 
      $$\bigcup_{i \in \mathbb{N}} A_{i} := \{ x \in X | \mbox{ il existe } i_{0} \mbox{ tel que } x \in A_{i_{0}}\}.$$
      $$\bigcap_{i \in \mathbb{N}} A_{i} := \{ x \in X | \mbox{ pour tout } i_{0} \mbox{ tel que } x \in A_{i_{0}}\}.$$
      $$A \backslash B := \{ x \in A \mbox{ et } x \not\in B\}.$$
      $$X\backslash A := \{ x \not\in A\}.$$
      $$ A \triangle B := (A \backslash  B) \cap (B \backslash A).$$ 
\end{defin}

\begin{prop}[Loi de Morgan]
     $$X\backslash \left(\bigcup_{i \in \mathbb{N}} A_{i}\right) = \bigcap_{i \in \mathbb{N}} (X \backslash A_{i}).$$
     $$X\backslash \left(\bigcap_{i \in \mathbb{N}} A_{i}\right) = \bigcup_{i \in \mathbb{N}} (X \backslash A_{i}).$$
\end{prop}

\begin{prop}[Distributivité]
     $$A \cap (B \cup C) = (A \cap B) \cup (A \cap C).$$
     $$A \cup (B \cap C) = (A \cup B) \cap (A \cup C).$$
\end{prop}
\end{document}